\subsection{有限表现模的同态}

设 E 为 A-模。若 I 是有向的预序集，且 \((\mathbf{G}_{i}, u_{ji})\) 是相对于 I 的 A-模归纳系，则典范映射 \(G_{i} \to \varinjlim G_{i}\) 诱导同态 \(\operatorname{Hom}_{\mathbf{A}}(\mathbf{E}, \mathbf{G}_{i}) \to \operatorname{Hom}_{\mathbf{A}}(\mathbf{E}, \varinjlim \mathbf{G}_{i})\) 从而得到一个称为典范的同态

\[
\varinjlim_ {i \in \mathrm{I}} \operatorname{Hom} _ {\mathrm{A}} (\mathrm{E}, \mathrm{G} _ {i}) \rightarrow \operatorname{Hom} _ {\mathrm{A}} (\mathrm{E}, \varinjlim_ {i \in \mathrm{I}} \mathrm{G} _ {i}).\tag{16}
\]

设 B 为另一个环，F 为 B-模，G 为 (A, B)-双模；我们在 II，p.～75 中定义了一个典范同态：

\[
\operatorname{Hom} _ {\mathbf {A}} (\mathrm{E}, \mathrm{G}) \otimes_ {\mathbf {B}} \mathrm{F} \rightarrow \operatorname{Hom} _ {\mathbf {A}} (\mathrm{E}, \mathrm{G} \otimes_ {\mathbf {B}} \mathrm{F}).\tag{17}
\]

命题 8. --- a) 若 A-模 E 是有限生成的（相应地有限表现的），则典范同态 (16) 是单射（相应地双射）。

\begin{enumerate}
\def\labelenumi{\alph{enumi})}
\setcounter{enumi}{1}
\tightlist
\item
  假设 B-模 F 是平坦的；若 A-模 E 是有限生成的（相应地有限表示的），则典范同态 (17) 是单射（相应地双射）。我们以 b) 为例证明，a) 的证明类似。将 A、B、F、G 视为固定的，并对每个右 A-模 E，令
\end{enumerate}

\[
\mathrm{T} (\mathrm{E}) = \operatorname{Hom} _ {\mathrm{A}} (\mathrm{E}, \mathrm{G}) \otimes_ {\mathrm{B}} \mathrm{F}, \quad \mathrm{T} ^ {\prime} (\mathrm{E}) = \operatorname{Hom} _ {\mathrm{A}} (\mathrm{E}, \mathrm{G} \otimes_ {\mathrm{B}} \mathrm{F})
\]

并记 \(\nu_{\mathbf{E}}\) 为同态 (17)；对每个右 A-模同态 \(v: \mathbf{E} \to \mathbf{E}'\) ，令 \(\mathbf{T}(v) = \operatorname{Hom}(v, 1_{\mathbf{G}}) \otimes 1_{\mathbf{F}}\) 和 \(\mathbf{T}'(v) = \operatorname{Hom}(v, 1_{\mathbf{G}} \otimes 1_{\mathbf{F}})\) 。

设 \(L_{1} \xrightarrow{v} L_{0} \xrightarrow{w} E \to 0\) 为 E 的一个表示；我们假设自由模 \(L_{0}\) （相应地自由模 \(L_{0}\) 和 \(L_{1}\) ）是有限生成的。图

\[
\begin{array}{c} 0 \to \mathrm{T(E)} \xrightarrow {\mathrm{T(w)}} \mathrm {T(L_ {0})} \xrightarrow {\mathrm{T(v)}} \mathrm {T(L_ {1})} \\ v _ {\mathrm{E}} \Big \downarrow \qquad \qquad \qquad \qquad \qquad \qquad \qquad \qquad \qquad \qquad \qquad \qquad \qquad \qquad \qquad \qquad \qquad \qquad \qquad \qquad \qquad \qquad \qquad \qquad \qquad \qquad \qquad \qquad \qquad \qquad \qquad \qquad \qquad \qquad 0 \to \mathrm {T^ {\prime} (E)} \xrightarrow [ \mathrm {T^ {\prime} (w)} ]{} \mathrm {T^ {\prime} (L_ {0})} \xrightarrow [ \mathrm {T^ {\prime} (v)} ]{} \mathrm {T^ {\prime} (L_ {1})} \end{array}\tag{18}
\]

是交换的，且其第二行是正合的（II，p.～36，定理 1）；此外，序列

\[
0 \rightarrow \operatorname{Hom} _ {\mathbf {A}} (\mathrm{E}, \mathrm{G}) \rightarrow \operatorname{Hom} _ {\mathbf {A}} (\mathrm{L} _ {0}, \mathrm{G}) \rightarrow \operatorname{Hom} _ {\mathbf {A}} (\mathrm{L} _ {1}, \mathrm{G})
\]

是正合的（同上引），且由于 F 是平坦的，(18) 的第一行也是正合序列（X，p.～8，定义 1）。既然如此，我们知道 \(v_{L_0}\) 是双射（相应地 \(v_{L_0}\) 和 \(v_{L_1}\) 是双射）（II，第 75 页，命题 2）。若仅假设 \(v_{L_0}\) 是双射，则由 (18) 可知 \(v_{L_0} \circ T(w) = T'(w) \circ v_E\) 是单射，因此 \(v_E\) 也是单射。若假设 \(v_{L_0}\) 和 \(v_{L_1}\) 二者都是双射，则由 X，第 6 页的推论 2 (ii) 可推出 \(v_E\) 是满射，而由于我们刚刚看到 \(v_E\) 是单射，故它是双射。

推论. --- 每个有限表现的平坦模都是投射的。

事实上，设 E 为有限表示的平坦 A-模。将 (b) 应用于情形 \(\mathbf{B} = \mathbf{A}\) , \(\mathbf{G} = {}_{\mathrm{s}}\mathbf{A}_{\mathrm{d}}\) , \(\mathbf{F} = \mathbf{E}\) ，我们看到典范同态

\[
\operatorname{Hom} _ {\mathbf {A}} (\mathbf {E}, \mathbf {A}) \otimes_ {\mathbf {A}} \mathbf {E} \rightarrow \operatorname{Hom} _ {\mathbf {A}} (\mathbf {E}, \mathbf {E})
\]

是满射。这意味着 E 是投射的（II，p.～77，注记 1）。

由前述推论和 X，第 10 页的命题 5，有限表示的平坦模与有限生成的投射模是同一类模。另一方面，存在有限生成的平坦模不是有限表示的，因此不是投射的（参见 X，第 170 页，习题 17；但另见 X，第 169 页，习题 13 和 14）。

